% ============================================================
%  Harald Høffding, "Pascal and Kierkegaard"
%  Revue de Métaphysique et de Morale, 30th year, no. 2
%  (April--June 1923), pp. 221--246 — the Pascal tercentenary number.
%  TRANSLATION (English), from transcription.tex.
%
%  NOTE ON THE SOURCE LANGUAGE: unlike everything else in this repo, the
%  original here is FRENCH, not Danish. Høffding wrote the article in French
%  for this journal. So this is French -> English, and TRANSLATION-PLAYBOOK.md
%  applies with the adjustments recorded under "Titles" below.
%
%  TITLES — policy fixed with the user, 2026-08-07:
%    - Kierkegaard's works: standard English titles.
%        Ou l'un ou l'autre                     -> Either/Or
%        Étapes de la route humaine             -> Stages on Life's Way
%        Postscriptum définitif non scientifique-> Concluding Unscientific Postscript
%        (L')Exercice dans le Christianisme     -> Practice in Christianity
%        (La) Maladie à la mort                 -> The Sickness unto Death
%        Le Moment                              -> The Moment
%    - Pascal's works: standard English, EXCEPT Pensées, which keeps its French
%      name because that is its standard name in English.
%        Lettres écrites à un provincial / Provinciales -> The Provincial Letters
%        Mémoire sur le Vide                    -> Memoir on the Vacuum
%        Comparaison des chrétiens des premiers temps avec ceux d'aujourd'hui
%              -> Comparison Between the Christians of Early Times and Those of Today
%        Apologie                               -> Apology
%        Discours sur les passions de l'amour   -> Discourse on the Passions of Love
%    - FRENCH SECONDARY LITERATURE STAYS IN FRENCH: Strowski, Pascal et son
%      temps; Sainte-Beuve, Port-Royal; the Revue itself; Œuvres. Translating
%      them would make the citations untraceable.
%    - Latin stays Latin: Augustine's De civitate Dei, Jansenius' Augustinus.
%    - The p. 232 footnote keeps Høffding's DANISH citation exactly as printed,
%      including the journal's misprint "Alsluttende" for "Afsluttende".
%
%  Fragment references "(fr. 144 Br.)" are Brunschvicg's numbering of the
%  Pensées and are kept verbatim.
% ============================================================
%
%  PAGE-JOINT AND PARAGRAPH REPAIR (2026 pass; fifth item swept, after
%        speculative-methode where the fault was found, evangelietroen-
%        theologien, om-personlig-sandhed and johannesclimacos):
%        Both files here were built by splicing page-range fragments into a
%        skeleton that separated them with blank lines. A blank line is a
%        paragraph break, so wherever the printed text ran on across a leaf the
%        splice invented a paragraph the article does not have.
%        Repaired in this file: 1 FALSE paragraph break (p. 239). No hyphen
%        joints and no missing breaks. The transcription alongside carried nine
%        false breaks and ten hyphens; see its header.
%        `paragraphs.py` in this directory performs the audit and takes either
%        file as its argument. Run it before any future publication.
%        ITS ONE HARD PROBLEM HERE IS THE RUNNING HEAD, which is full-measure
%        and lies in BOTH directions: on even pages the Revue sets the folio out
%        in the left margin, so the head begins ~340 px LEFT of the block and the
%        page reads as a -337 px hanging indent; on odd pages the head is
%        CENTRED, so it begins ~100 px RIGHT of the block and is
%        indistinguishable from a paragraph indent — printed p. 245 was reported
%        as +97 px on that mistake and in fact opens flush. Two guards are needed
%        together: the line must sit in the text block AND not be followed by the
%        extra white space that separates a head from the body. With both,
%        continuations measure -19 to +20 px and the one paragraph opening +77.
% ============================================================
\documentclass[12pt,a4paper]{article}
\usepackage[utf8]{inputenc}
\usepackage[T1]{fontenc}
\usepackage{libertinus}
\usepackage{libertinust1math}
\usepackage{textalpha}
\usepackage[danish,main=english]{babel}   % Danish for the one quoted citation
\usepackage[protrusion=true,expansion=true]{microtype}
\usepackage{setspace}
\usepackage[margin=1.2in]{geometry}
\usepackage{enumitem}
\usepackage{fancyhdr}
\setlength{\headheight}{28pt}
\addtolength{\topmargin}{-13.5pt}

\newcommand{\dk}[1]{\foreignlanguage{danish}{#1}}

\usepackage{hyperref}
\hypersetup{
  pdftitle={Pascal and Kierkegaard},
  pdfauthor={Harald Høffding},
  hidelinks
}

\pagestyle{fancy}
\fancyhf{}
\fancyhead[LE,RO]{\thepage}
\fancyhead[RE]{\textit{Revue de Métaphysique et de Morale}}
\fancyhead[LO]{\textit{Pascal and Kierkegaard}}

\setstretch{1.3}

\begin{document}

\begin{center}
  {\Large PASCAL AND KIERKEGAARD}\\[1ex]
  {\normalsize Harald Høffding}\\[0.5ex]
  {\footnotesize\itshape Revue de Métaphysique et de Morale 30:2
   (April--June 1923), pp. 221--246}
\end{center}

\bigskip

% The article opens with an unnumbered introduction, pp. 221--224; section I
% begins on p. 225. Structure mirrors transcription.tex 1:1.

% --- p. 221 ---
Two centuries apart, two thinkers --- both of them bound to Christianity to
the very depths of their fervent souls --- met in bringing suit against their
religion as the ages had deformed it. Each denounces the incompatibility that
sets the spirit animating Christianity at its beginnings against its condition
in modern times. This antagonism dates from the day the Church attempted to
enter into intimate relations with the ``world,'' the ``age'' --- in other
words, with a culture developed on purely human foundations. In both cases the
reckoning with the Church was undertaken with such passion, and pressed with
such logic, that there was matter here for a great intellectual and moral
revolution; but --- and this is precisely the tragic point in the destinies at
issue --- scarcely had the two thinkers uttered their last word than they
succumbed, each consumed by the spiritual struggle he had had to sustain with
himself and with his contemporaries. They died young, Pascal at 39 (1662) and
Kierkegaard at 42 (1855), leaving untouched the great problem they had
formulated anew with unprecedented sharpness. But in the world of the spirit
heroic energy is never spent in vain, and, taken up with such urgency and such
fire, the question arose whether or not the religious traditions of the
Christian Churches could continue to support the honest disquiet and the
sincere searching of the friends of truth.

A vain question! from the camp of the Churches no answer has come.
% ^ The original prints an exclamation mark followed by a LOWER-CASE "du";
%   the transcription keeps it, and the translation mirrors the run-on.
There they flatter themselves that they persist in the practice they have
followed, by disregarding the recriminations. They lean on the fact --- an
incontestable one, moreover --- that up to now no other organization has
appeared capable at once of spreading through the world the notion of the
seriousness of human existence
% --- p. 222 ---
and of bringing men the relief and the comfort that its hardness demands. So
the problem remains in suspense. It is told of Scipio Africanus that, accused
before the assembly of irregularities in his accounts, he cried out: ``Romans,
on this same day I conquered Hannibal at Zama; let us go up to the Capitol and
render thanks to the gods for it!'' The people followed him, and there was no
further question of rendering accounts. The doubt has kept all its force since
the two thinkers, and historical studies have even deepened its foundation.
% ^ The original prints "appronfondi" for "approfondi" (a compositor's error,
%   kept in the transcription). Rendered here by the intended sense.
Have the Churches the right to believe in a living continuity between
themselves and the great spiritual movement that set out more than nineteen
centuries ago from a corner of the Roman empire, and whose governing idea
commanded men to live thenceforth in extreme expectation of the imminent and
supernatural advent of a new kingdom --- an expectation that reduced to a
trifling value the human frames of existence, the family and the State, art
and science, when it did not see in them so many obstacles to the kingdom of
God?

It is not, moreover, this suit against the Churches alone, whatever gravity it
may have had, that makes us remember Pascal and Kierkegaard. Before concluding,
their minds developed ideas of great importance: here, as often elsewhere in
the domain of thought, the road travelled has its full value independently of
the goal attained. Along the way they cast light upon the nature and the
conditions of existence of man's spiritual life, reaping without stint
psychological findings of high price. They made contributions to what we today
call the theory of knowledge, by bringing thought into play at the outermost
limit of its domain. They showed that even where feeling and passion carry the
word, it is still the idea that determines --- like a guiding star --- the
direction of psychological movement. They are the champions of Thought, and
their struggle remains interesting whether or not it ended in victory. They
followed different roads. Their individualities appear, within and without,
under different aspects, each presenting constitutive traits of its own, and
the historical settings in which their thoughts were sharpened were likewise
very dissimilar, while yet offering certain analogies between them.

In their style, already, they differ. We must recall here that Pascal's
\emph{Pensées}, which constitute his principal work, do not present themselves
to us in the definitive form of a work
% --- p. 223 ---
brought to completion by its author. They are fragments, notes scribbled by
Pascal in the respites his illness allowed him, often interrupted by the onset
of an attack. Concise, compressed, vibrating with an emotion that in an
alternating rhythm makes the thoughts spring forth or springs forth from them,
these fragments have a force and a clarity without equal, and there is no
reason to believe that these qualities would have been lost in a definitive
redaction. Kierkegaard's style has neither that powerful stamp nor that
density. For him --- at least in the first phase of his literary production ---
the work of composition was a cure: so long as he was writing, melancholy had
no hold on his soul. And so he took his time. Besides, he loved language for
its own sake; he played with it without always noticing that the reader, for
his part, might be a little more pressed, impatient to take in as a whole the
sequence of ideas, to see what all this might finally come to. Lyrical
digressions, descriptions of nature, episodes borrowed from everyday life come
to interlace themselves with the chain of the argument. The passion that lies
at the base of his production, as of Pascal's, makes itself felt in the form of
a continuous tremor rather than of emotional crises. Only in the last period of
his production does that tremor turn into an earthquake.

If, in this comparative study, I happen to speak at somewhat greater length of
the Danish thinker, it is because I am addressing French readers. In France the
study of Pascal has given rise to a production so abundant, so solid and so
well documented that to claim to add anything would truly be to carry owls to
Athens. Let me be permitted, in this connection, to express my admiration for
Sainte-Beuve's \emph{Port-Royal}, which I have just reread at a distance of
nearly 60 years and which, after so many works published since, keeps all its
freshness and all its penetration. I must also pay homage to that splendid
edition of Pascal's \emph{Œuvres} which we owe to M. Léon Brunschvicg and his
collaborators MM. Pierre Boutroux and Félix Gazier, and whose
\emph{Introduction} and annotations furnish elucidations of great value. I am
grateful above all to M. Brunschvicg for having set out, in an
\emph{Introduction} to the \emph{Pensées}, in an extremely interesting manner,
the logical sequence of ideas that would probably have been Pascal's in the
\emph{Apology} he proposed to write. In the lines that follow
% --- p. 224 ---
I shall refer, for citations of the \emph{Pensées}, to the order proposed by
M. Brunschvicg. --- As for Kierkegaard's works, a second edition of his
complete works is at present in course of publication. A German edition has
appeared (Diederichs, Jena). They comprise fourteen volumes. But, in addition,
the integral publication of Kierkegaard's \emph{Journal} and private papers has
been begun (by P.-A. Heiberg and V. Kuhr). Up to now eleven volumes have
appeared, and there remains the matter of seven volumes more. We are therefore
in the presence of a very voluminous production. In French, a fine study of
Kierkegaard by M. Delacroix was published in 1900 in the present \emph{Revue},
which also gave, in 1913, an address delivered by the writer of these lines at
the University of Copenhagen to celebrate the centenary of the Danish
philosopher (5 May 1913). A monograph on Kierkegaard as philosopher, which I
wrote in 1892, has been translated into German and is about to be translated
into Spanish.
% ^ The original prints "allemant" for "allemand"; rendered by the sense.

The hypothesis of an influence exercised by Pascal on Kierkegaard appears less
than probable. It is true that Pascal's works, in German translation, figure in
the inventory of Kierkegaard's library; but it is only late (in the autumn and
winter of 1850\footnote{On information I owe to the kindness of M. P.-A.
Heiberg, one of the editors of the \emph{Journal}.}) that the latter makes
mention of Pascal in his
% ^ Footnote 1 of the article. In the original the mark is a superscript arabic
%   numeral followed by a closing parenthesis on the baseline ("1850 ¹)"), and
%   the note is set at the foot in small type with no rule above it.
\emph{Journal} --- that is to say, at a moment when he had already produced the
writings [\emph{The Sickness unto Death} (1849) and \emph{Practice in
Christianity} (1850)] that show the greatest affinity with Pascal's ideas. He
read, thereafter, various works dealing with Pascal, notably Reuchlin's, on
Port-Royal, and he noticed the difference between the older editions of the
\emph{Pensées} and Faugère's, in which the most revolutionary notes are
admitted. (The German translations represented in Kierkegaard's library had all
been made from the earliest editions.)

In the comparison of the two thinkers that I am about to undertake, I shall
stop successively at their temperament and their character, at their
intellectual antecedents, at their relation to Christianity, and at the
question of what paths they left open at the moment when they fell silent.

% --- p. 225 ---
\section*{I. --- \textsc{Their temperament and their character.}}

Pascal's physical and moral temperament has been defined in concordant fashion
by his two sisters (Jacqueline's letters, M\textsuperscript{me} Perier's
biography). He was of a boiling humour, and the vivacity of his mind did not go
without an impatience that made him at times difficult to satisfy; his elder
sister marvelled that with so fiery a temperament he had been preserved from
the excesses of youth. A thought that seems to me altogether significant in
this regard is that of fragment 139 (Br.), according to which all the unhappiness
of men would come from a single thing, which is ``not knowing how to remain at
rest.'' In the curious \emph{Discourse on the Passions of Love}, which is
attributed to Pascal, it is the beauty of the impetuous transport of love that
is magnified, and passion in general is characterized by the precipitation of
all one's thoughts in a single direction. Here is announced that passion of
faith which was to take hold of him after his definitive conversion. His nature
kept its primitive character through all the phases of his religious evolution,
and although he had to pass through a double conversion before arriving at the
term of his development, Pascal was never a man without religion. Throughout his
youth his soul was dominated by a feeling of religious gravity; but alongside
it an intellectual passion, precociously active, engaged him in researches that
made an epoch in the history of mathematics and physics; nor did this prevent
him for a time from going into society and seeking there the intellectual and
literary agreeableness of the company of well-bred people. It is to this worldly
period that we must refer the adventure which probably inspired the
\emph{Discourse on the Passions of Love}. According to M. Strowski
(\emph{Pascal et son temps}, II, p. 278) the object of his passion appears to
have been a beautiful lady of high birth, of an adventurous and passionate
humour, but who gave him no matter for repentance. He was nonetheless, later
on, gravely to impute it to himself as a sin to have had his soul filled with
anything other than God. For in Pascal's eyes there was here a mortal antinomy,
an antagonism equivalent to the conflict of life and death. He was of those
natures who need a total rupture with the past in order to reach their
definitive position. His ``boiling humour'' required it. The catastrophe came
in that night of 23 November 1654 when, in an ecstasy, he
% --- p. 226 ---
understood that the living God demands everything and that a plenary
renunciation must follow. His life up to that day seemed to him spent in the
void; now he exulted at having found again the source of life.

Having passed through such psychological states, Pascal could only be indignant
at the attempts made by the Fathers of the Society of Jesus to adapt the high
demands of the Christian ideal to the needs of people of the world. Impossible
to ``join God to the world'' (fr. 935 Br.). The world does not repel religion,
but it wants it wholly indulgent and gentle (fr. 956 Br.). In his \emph{Provincial
Letters}, as in the very last of the \emph{Pensées}, he speaks in
ever harsher terms of the ways of corruption into which the Church has sunk.
``If she absolves or binds without God, she is no longer the Church'' (fr. 870
Br.).

In the course of the polemic thus engaged, at the very time when he was
publishing the \emph{Provincial Letters}, it happened that a young relative of
Pascal's was cured of a disease of the eyes that the physicians had nevertheless
declared incurable: there had been applied to the diseased eye a thorn from the
Saviour's crown, which formed part of the relics of the church of Port-Royal.
This miracle confirmed him in his conviction of having chosen the true. The hour
of miracles had therefore not passed, as his adversaries maintained. And if the
authorities of the Church deny the truth, God himself must intervene (fr. 832,
839 Br.). And so Pascal appeals with confidence, from Rome, which condemns him,
to the tribunal of Jesus Christ (fr. 920 Br.).

In his private life Pascal realized complete renunciation of himself. Without
entering a convent, he lived his last years in the most rigorous asceticism,
wholly given over to works of charity, and anxiously forbidding his own to
attach themselves to his person and so to defraud God of that total surrender
of one's being which is owed to him. The long illness that was consuming him he
took as a relief, although it prevented him from finishing that apology for
Christianity whose sketch is contained in his \emph{Pensées}. For the Christian
the normal state is a morbid state of the soul, in which it is less difficult
for him to maintain himself when he suffers in his body as well.

\bigskip
% ^ The original sets a blank line here --- visibly more than a paragraph
%   break, less than a section break --- dividing the Pascal portrait from the
%   Kierkegaard portrait. No ornament, no rule, no numeral.

It is another moral portrait that emerges for us from Kierkegaard. His
development has more continuity, while yet offering decisive turning points. He
does not belong to the category of the ``twice born.'' Neither ecstasy nor
miracle plays any part in his life.

% --- p. 227 ---
Like Pascal, he was brought up in a Christianity full of seriousness, but
during his youth he occupied himself with poetry and philosophy. It was the
time of Romanticism, and the idea of a harmony embracing art, science and
religion dominated men's minds. Romantic ideology had just been reduced to a
system by Hégel,
% ^ "Hégel" with an acute accent, as printed; ten lines further down the same
%   page the original sets "Hegel" unaccented. The inconsistency is the
%   journal's and is left standing in both files.
and Kierkegaard took an interest in it for a time. Since the antinomies all end
by resolving themselves into harmony, one might well at first give free rein to
one's sentimental reverie and one's imagination; do they not form part, indeed,
of that rich and changing world whose enigma could, in the last analysis, be
resolved by the essential truths of Christianity? But early on Kierkegaard
perceived that at this game one does not escape an irreducible duality: it is
in this sense that we see him note down at that time in his \emph{Journal}, as
a motto, Goethe's line about Gretchen: ``Halb Kinderspiel, halb Gott im
Herzen!'' (Children's games and God dwell side by side in her heart).
% ^ the German is set in roman inside the guillemets, not in italic.
He even had at that time a period in which he let himself be drawn into a life
of dissipations which was afterwards to feed with long remorse that melancholy
inherent in his nature and at moments so powerful that he felt himself on the
edge of madness.

He set himself thenceforth the mission of fusing into a personal harmony the
contradictory impulses of his being --- unconcern and gravity, frivolity and
melancholy. Personal experience then showed him that it was far more difficult
in practice to attain that ``higher unity'' which Hegel's system realizes
\textit{de plano}. This inner labour made a solitary of him. He felt isolated
and ill understood. His melancholy bent forbade him, he thought, to enter into
bonds implying a sincere and joyous self-surrender. He broke off his engagement
to a girl of gay and playful humour, and judged himself incapable of exercising
any function whatever. The contrast was too glaring, in his eyes, between the
light ease with which others took existence and the struggle in which he
wrestled against the darkening forces of his life. He told himself that in the
Middle Ages his entry into a monastery would have been clearly indicated. This
obscure and patient elaboration of his personal life he described under this
poetic image: ``I listen for my inward music, for the joyous calls of its song
and for its grave bass organ notes. And it is no small task to coordinate them
when one is not an organist but a man who confines himself, for want of larger
% --- p. 228 ---
demands upon life, to the simple wish to want to know himself.'' --- How far we
are, then, from the time when ``God'' and ``children's games'' dwelt at peace
side by side in his heart. Their antagonism constitutes a problem whose
solution demands the whole energy of a life. And in fact the effort to
reconcile the irreconcilable is characteristic of Kierkegaard's whole life. He
does not go from one extreme to the other: the contraries are there, confronting
one another, from the first; only their opposition grows sharper as the
development of his individuality advances, and the effort to overcome them
increases in proportion. It is on the scale of his own experience that
Kierkegaard was able to measure the intensity of effort of the personal life.

One means offered itself to him for holding melancholy in check and easing the
inner tension: he had only to give way to his inclination to literary
production. It was a deliverance for him to give himself over to literary
production.
% ^ The phrase "à la production littéraire" is set TWICE in two consecutive
%   sentences in the original. Almost certainly a compositor's dittography, but
%   it is what the page carries, so the repetition is reproduced here too.
Like the Princess of the \emph{Thousand and One Nights} who told tales to save
her life, he saved his own, he says, by writing. To the writer's first period
(1843--1846) go back some of his best-known writings: \emph{Either/Or} and
\emph{Stages on Life's Way}. They turn on the very object of his inward
struggles: the interior elaboration of personality, the conquest of a solid
core about which the other elements of the soul will gravitate. By this his
first works could serve besides to stimulate the enervated generation of
Romanticism and speculative Philosophy, whom the epigones were leading to blunt
the antagonisms of man and thereby to debase his spiritual life. Kierkegaard
could not claim that he had had such an influence in view at the beginning of
his career. Honestly he recognized that there had been in him a tension that
called for an outlet. Unlike Pascal, who deliberately proposed in his
\emph{Provincial Letters} to combat a collapse of morality, Kierkegaard came
little by little, and without really having willed it, to work in the same
direction. He was besides not yet at that moment as far advanced as Pascal in
his religious development.

A new period in Kierkegaard's career (1847--1855) opened precisely because his
inner life took on a decidedly more and more religious character. He writes in
his \emph{Journal} (1847): ``I now feel the need of an ever deeper
understanding of myself by drawing ever
% --- p. 229 ---
nearer to God. Something is taking shape in me, I know not what, that announces
a metamorphosis\ldots{} I must therefore keep still.'' One will remark the
contrast with Pascal's ecstasy: here, the almost unconscious emergence of a new
mode of existence, before which the subject keeps a passive attitude. The
following year he notes: ``At present I possess faith in the intimate acceptation
of the term.'' Christianity was taking on in his eyes a character of reality
that until then he had not known in it. This first period of his career ended,
he believed for a time that his literary activity too had come to an end; the
idea smiled on him of withdrawing to some remote corner in the country. But he
did not carry it out: the increased force of his religious consciousness
awakened a severer criticism of existing Christendom. In his first literary
period Kierkegaard had branded the laxity into which the personal conception of
life had fallen; the Christian now discovered in the Church a readiness to
compromise, a spirit of accommodation of Christianity to the tastes of worldly
people, which relegated to the background the ideal of the first Christians. It
would be, according to him, virtually to abolish Christianity to make of it a
religion wholly of gentleness and consolation. ``And we see living, in present-day
Christendom, a generation spoilt, proud, and yet cowardly, arrogant but without
spring, which lets itself be administered from time to time these good
consolatory principles, without even knowing exactly whether it will make use of
them when life smiles on it, and which takes offence at them in hours of
distress when it appears that at bottom their indulgence gives way.'' We borrow
this citation from \emph{Practice in Christianity}, which characterizes, with
\emph{The Sickness unto Death}, this second period of Kierkegaard's production.

The official Church did not take up the observation of a glaring conflict
between the obligation harshly imposed by primitive Christianity to pursue
nothing but our one thing needful and, on the other hand, the idyll organized by
modern Christendom under cover of the dogma of Redemption. It was then that
Kierkegaard began his passionate war against the existing Church, the most
violent of all the struggles our Danish intellectual history has known.
Kierkegaard's last word in it was: ``The Christianity of the New Testament does
not exist.'' In the thick of the fight illness stopped him, and death.

% --- p. 230 ---
\section*{II. --- \textsc{Intellectual predispositions.}}

Pascal pursued his studies in mathematics and physics and won a fine name in
them. It is to these studies that he owed his conception of scientific
knowledge. A system like Descartes's, tending to reduce knowledge of nature to
certain fundamental propositions and so to prepare rational foundations for the
work of human reason --- such a system seemed to him condemned in advance. He
has justly been given the name of experimental positivist. From his ``worldly''
period he had understood that, given the multiplicity, the complexity and the
variety of psychical phenomena, it is not from the ``spirit of geometry'' ---
that is, from reasoning proceeding by purely mathematical deductions --- that
one must ask for an understanding of what goes on in the human soul. In
psychology what matters above all is to refer the multiple phenomena to certain
types, to identify oneself in thought with each of these types, and to seek to
know their motives and their preoccupations. Here mathematical reasoning would
be found wanting; only the ``spirit of finesse,'' a subtle sense for types and
for nuances, will allow us to disentangle the typical elements and to find them
again in their varied combinations. M\textsuperscript{me} Perier admired in her
brother precisely this aptitude for taking account of the mental data of his
interlocutors. He knew, she thinks, ``all the springs of the human heart.'' This
faculty had been developed in him by frequenting people of the world; it proved
useful later, when it was a question for him of winning souls for the cause of
the true Religion. The proofs furnished by speculation are of little efficacy.
Conclusive or not, their action is confined to the moments when we are under
the influence of intellectual movements; afterwards they are easily forgotten.
The good tactic is to take men at the point where they have stopped and to show
them how to go forward by their own resources: nothing convinces men better
than the reasons they draw from their own stock. They must be helped to see the
curious mixture of greatness and of misery that their nature is. Thought may
well carry them up into its high regions, but concupiscence and the passions
will always bring them back to the mire below. Man must be brought to
understand what an ``incomprehensible monster'' he makes, and
% --- p. 231 ---
only the true Religion will furnish him the key to his nature together with the
ways and means of resolving the inward divisions that tear him. Pascal is
concerned above all that man should come to feel himself alone in an infinite
world from which no voice of compassion or encouragement issues. The trouble,
the disarray, the anguish that Pascal felt in the face of the great enigmas made
him take the only way that opened, the issue of Christianity. No doubt
Christianity contains mysteries alongside the revelations it brings to men. But
if God had left us ignorant of nothing in the mystery of our life in the
universe, would he not at the same stroke have suppressed in the soul that
effort, that combat, that courage to dare and to stake everything, which alone
allow it to merit eternity?

Pascal reduces to two principal types the men who attempt to resolve, by
natural lights, the great enigmas of existence. The one sort attach themselves
to thinking reason and to duty, which are properly the greatness and the solid
point in man, and the rest they hold to be negligible. Their classic
representative is Epictetus the Stoic. Minds of the other type accommodate
themselves to the fluctuations of life; they practise the \emph{carpe diem}
and keep besides to received opinions, being well aware that everything is
subject to vicissitude. They repose with delight in the certainty of being sure
of nothing. Montaigne is the paragon of the conscious unconcerned.

The study of these two types was for Pascal a preparation for the hand-to-hand
struggle in which he proposed, in his \emph{Pensées}, to seize and shake the
unbeliever, in order to put him in the critical state in which he would perceive
only one single door of salvation --- the one Pascal himself had found. It was
a disappointment to him to see how little the study of human nature was
appreciated, the investigation of its types and of their practical utility.
``When I began the study of man, I saw that the abstract sciences are not proper
to man, and that I was straying further from my condition in entering into them
than others were in being ignorant of them.'' He had hoped, then, to find more
numerous companions in the study of man; but here, contrary to his expectation,
he found even more detachment than in the case of geometry (fr. 144 Br.). He
himself threw himself headlong into the new study, growing more and more
indifferent to his earlier scientific labours. In a letter to the celebrated
mathematician Fermat he writes towards the end of his life (10 August 1660)
that, while still
% --- p. 232 ---
regarding geometry as the highest domain in which the faculties of the mind can
be exercised, he now looks on it as altogether useless, especially since he has
undertaken studies so different from his earlier work that the latter has almost
vanished from his memory.

He remains convinced that man's greatness and titles of nobility reside in
thought, and that it is to thought that we must attach ourselves. Compared with
the great masses that surround us we occupy but little space, and our power is
negligible beside theirs. We are reeds, but thinking reeds, and though the
external masses crush us, they will never attain the dignity that thought
confers (fr. 346--348). But this thought, so exalted by Pascal, was reduced more
and more, for him, to being no more than a servant of theological explanation.
Just as, during his ``worldly'' period, he had described the passion of love as
consisting in a concentration, a precipitation of all one's thoughts towards a
single end, so now that religion had become his dominant passion, he admitted
only thoughts oriented in that single direction.

In Kierkegaard's case the intellectual data are quite different. Grown up among
the epigones of Romanticism, he saw for some years, in Hegelian metaphysics,
the philosophical ideal, whereas in Pascal's eyes ideal thought was represented
by geometry. In a system like Hegel's, where an inner necessity binds all the
parts to one another, there is nothing that has not its logically marked place.
Poetry, Religion, Science were therefore for Kierkegaard only the diverse or
successive forms of one and the same divine truth. Such was the dream of his
youth, from which he finally freed himself --- following his teacher, Professor
Paul Moller,
% ^ As printed: "Paul Moller", for Poul Møller. Høffding's own French form of
%   his teacher's name, without the ø. Kept in both files.
the most Danish of all Danish writers. Kierkegaard's merit, in the philosophical
order, is to have been one of the first to criticize speculative philosophy from
the standpoint of the theory of knowledge (\emph{Concluding Unscientific
Postscript}, 1846)\footnote{\dk{\emph{Alsluttende uvidenskabelig Efterskrift.}}
Copenhagen, 1846.}%
% ^ Høffding's Danish citation is kept exactly as the journal prints it,
%   including its misprint "Alsluttende" for "Afsluttende" (verified at 600
%   dpi), so that transcription and translation agree.
.
He shows that even if an absolute knowledge existed, each of us would
nevertheless be obliged to begin, at some determinate point
% --- p. 233 ---
taken within the domain of experience, the reasoning that was to raise us to
that higher stage --- and why, then, at one point rather than at another? And
this point of departure, chosen in the real world, will one be able in turn to
deduce it from the absolute knowledge that is to result from it? And will there
not always subsist other data besides that experimental point of departure at
the very base of absolute knowledge, and what certainty will one have of having
gone back to the remotest of them? Kierkegaard was led to see that of all
philosophical problems that of knowledge is the most central. He insists on the
fact that we live in time. We live forwards, but we understand backwards. Hence
the impossibility of a terminal conclusion; and Kierkegaard declares himself in
agreement with Lessing in saying that for man there is nothing beyond the search
for truth.

Like Pascal, Kierkegaard is of the opinion that the personal life needs
thoughts, but thoughts applied to what shores up our existence. The thought of
God is only a desperate expedient and offers no object to pure intelligence.
Kierkegaard establishes a distinction between ``essential'' and ``non-essential''
knowledge, classing among non-essential knowledge: the natural sciences,
philology, history and pure philosophy --- all of them thoughts without
employment for us when we take up a position in the capital problem of life and
death. Impossible to conceive life in a purely intellectual fashion, unless one
makes an abstraction of oneself. To appraise the various attitudes men take up
towards existence, Kierkegaard uses a procedure recalling Pascal's: he groups
them in a series of types which he undertakes to characterize well. Let us
retain here only those whose portrait he drew in his most popular writings
(\emph{Either/Or} and \emph{Stages on Life's Way}). The ``aesthete'' wants
above all to enjoy the passing moment. His art of living is to keep enough
elasticity of soul to seize the successive novelty of the instants. It is
important to break off at the opportune moment, so as not to be held in
immobile bonds. The method of rotation of crops has something to be said for it.
The aesthete is like a tangent to the circumference (more exactly: to the
circumferences) of life. At a single point there is, momentarily, a contact,
which is at once followed by a contact at another point. Hence: loves that
change and no marriage; comradeships and
% --- p. 234 ---
no friendships; improvised occupations and no professional duties; intellectual
enjoyments and no science! If the ``aesthetic'' life is a tangent, the
``ethical'' life describes a circumference about a centre and finds itself
determined by relation to that centre. It is the type of fidelity. Its life is a
life in common, a regulated life, a life of incessant work. Above these two
types rises that of the religious life, whose effort and orientation refer to an
end situated beyond all moments and all duration, incommensurable with all other
human ends and even, at its culminating point, in Christianity, in manifest
contradiction with everything that, from a purely human point of view, might be
taken as the end and the directive of existence. The more the character of the
religious life is differentiated, the less will it allow man to develop in a
natural environment: he will be like a fish on land. Finally, Eternity will not
remain (as with Plato and with Spinoza) the immense, ever-present ground upon
which the life measured in time unfolds: it will irrupt into duration, the
divine arising in human form at a determinate point of time and space. Then this
revelation of men's attitude towards it will be the decisive factor for
appraising life and its content.

It is a comparative philosophy of life that Kierkegaard gives us here. It calls
to mind the characters of Montaigne and of Epictetus,
% ^ As printed the original has "Epictète" without the acute here, though the
%   same name is set "Épictète" on p. 231. Both recorded, not normalised.
described by Pascal, and that of the Christian whom he opposes to them both.
And just as Pascal saw in illness the state natural to Christians, so
Kierkegaard, for his part, arrives at the result that the highest type of life,
that of Christians, involves tension and sufferings owing to the lacerating
antagonism of the contradictories between which life must be lived. And the
superiority of this last type over the others comes precisely from the greater
trials it imposes on that power of synthesis which every personal life demands.

According to Kierkegaard, it is by leaps that one passes, as from stage to
stage, from one life to another, and the principal leap is the one that carries
a soul plunged in despair to life in Christ. But even at the most decisive
point this leap never takes on, in Kierkegaard, the character of ecstasy. And in
a certain measure continuity subsists, since our synthetic energy remains the
scale common to all types of personal life. Besides,
% --- p. 235 ---
Pascal at bottom employs the same measure; he draws it from the relation between
the more or less rich stock an existence offers and man's energy in
concentrating and ordering its riches. It is indeed the only measure one can
employ with respect to differences, of individuals and of epochs, in their
manner of understanding life; it imposes itself even on those who do not place
the final end of existence where Pascal and Kierkegaard put it. An objection
comes to us at once against this use they make of their principle of
measurement. They consider exclusively the inner, subjective resistance that has
to be overcome, and consequently the inner labour by which the unity and the
concentration of the soul are obtained. But the question arises what the
importance of this inner labour is for the human groups in which the individual
lives, for the society of which he forms part, perhaps even for humanity
entire. This is a point of view one would never get Pascal or Kierkegaard to
adopt. According to them, there is a whole abyss between the type of life they
regard as the highest and the other lives consecrated to the work of human
civilization.

\section*{III. --- \textsc{The Christian problem.}}

Pascal's life and work were accomplished within Catholicism; Kierkegaard
developed in a Protestant milieu. Hence, between them, multiple differences
showing themselves in their religious formation as well as in the positions in
which they definitively established themselves. But from the point of view of
the history of religion it is very interesting to see that, having arrived at
the term of their effort to discover the central point of their religion, they
halt before one and the same problem, which might be formulated thus: Is it
possible to remain in the spirit and the practice of primitive Christianity
while assimilating and serving a culture, moral and material, which not only
does not issue from Christianity but even preceded it, and has continued to
develop outside it according to its own particular laws?
% ^ A single damaged sort stands immediately before "s'assimilant" in the
%   original and is not legible as any letter at 600 dpi; the sense wants
%   "tout en se l'assimilant". Rendered here by the sense.
In effect, the problem arises for Catholics as for Protestants. Before going
further, let us consider what gave rise to it.

The first Christian generations lived in expectation of a near return of Christ,
which was to close the destinies of the world by a universal judgment. At the
outset this expectation constituted the belief of Christians; it regulated their
life and their conception
% --- p. 236 ---
of life. In the perspective of an imminent catastrophe out of which a new heaven
and a new earth would be born, purely human or social duties lost their urgency.
That is why there is no question, in the New Testament, of art, of science, of
public life; family life itself was admitted only under certain reserves. The
morality of the New Testament had to be what M. Albert Schweitzer has aptly
named an ``interim ethic.'' All efforts, all projects converged on a single end
set in a neighbouring future. For this primitive Christendom the advent of the
kingdom of God, alluded to in the Lord's Prayer, took on a very precise sense:
it was not a matter of vaguely asking thereby for the obtaining of spiritual
goods. Among the fervent, expectation was doubled by an enthusiasm often raised
to the pitch of ecstasy, which made all sacrifices easy. The things of men's
daily life shrank before the grandeur of those that were awaited.

As against a belief like that of the primitive Christians, Catholicism and
Protestantism appear as long-term organizations based on the chance of a
prolonged evolution of religious life and, in general, of human life on earth.
Less and less has one prepared for an abrupt departure. While striving not to
depart from the expectation of the first centuries, its realization has
nevertheless been placed in an ever more remote future, and the subjection in
which it had held souls went on diminishing. Little by little men appropriated
--- after a suitable resistance, which was often harsh enough --- not only the
culture transmitted by classical antiquity but also the new culture that was
developing independently of Christianity: a fact which does not prevent the two
Churches from asserting themselves in continuity with the belief of the first
Christian generations and from maintaining, each in its own manner, the
pretension of having kept the old tradition intact through the new times.

The most grandiose solution of the problem is that of Catholicism. It is founded
on the distinction of degrees in perfection. Monks and nuns renounce life in the
world in order to be able to consecrate themselves entirely to the beyond, to
the one thing needful. It is through them that the filiation of original
Christianity subsists. To the laity falls the task of continuing human life under
% --- p. 237 ---
the influence of the directors of conscience through whom the Church exercises
its action of control over souls. Saint Augustine (\emph{De civitate Dei}, XX,
9) already saw, in the prelates of the Church, a partial realization of the
prophecy of the Apocalypse, speaking of those who occupy the tribunal in the
kingdom of a thousand years.

Pascal understood very well that in the eyes of the Church the estrangement of
original Christianity from all forms of culture must never cease to prevail in a
certain practice. Without entering a convent, his life in his last years was
nonetheless a monk's. In his eyes the times that separate us from the return of
Christ must be nothing but ardent expectation and suffering. The scene of
Gethsemane represents well, according to him, the permanent situation of the
true Christian. We must live as active contemporaries of the Saviour's
sufferings. Once only, at Gethsemane, Jesus in distress asked help and comfort
of men, and his apostles failed him. But the duty in which they then failed
falls to the Christians of successive ages. ``Jesus will be in agony until the
end of the world; we must not sleep during that time'' (fr. 583 Br.). Therefore
it is not towards the victorious Christ that we must turn our gaze, as the
Church most often does, but towards the suffering Jesus. The garden of Olives is
the counterpart of the garden of Eden. Fallen man must make of Gethsemane the
permanent dwelling of his thoughts; at that price only does he keep the hope of
seeing the gates of Eden open before him.

The direction of souls as the Jesuit fathers practised it appeared to Pascal as
a criminal accommodation of life according to Christ to life according to the
world; and in \emph{The Provincial Letters} he raised a lively protest against
the abuses of the Jesuits. Save for errors of detail, he is right from the
rigorously Christian point of view. From the New Testament to the casuistry of
the Jesuits the ethic does not remain identical, and even from a purely human
point of view one might raise doubts touching the results of this casuistry.
Alongside his campaign against the religion of gentleness which had, according
to him, been substituted for the meditation of the sufferings of Gethsemane,
Pascal devoted a memoir to the \emph{Comparison Between the Christians of Early
Times and Those of Today}.

If there is so wide a gap between the two categories of Christians, it is
because at present baptism precedes instruction in the mysteries of
% --- p. 238 ---
religion, whereas originally no one had access to Christianity without having
been well instructed in its doctrine and without having renounced the world. For
then the distinction imposed itself between the ``world'' and the Church,
whereas today the Church has admitted the world into its bosom by granting
baptism to children. The institution of infant baptism, Pascal adds, was made
with good intention: it was feared that children who died unbaptized would be
deprived of salvation; but what was done for the good of children has turned to
``the ruin of adults!''

Pascal's rupture with the Church was not brought about by the divergence
observed on this single point. As we have already seen, he had for rising
against the authority of the Church reasons drawn not only from the practice of
the Church but further from its dogmas; it was also a question of the
establishment of facts. In agreement with his Jansenist friends, Pascal
maintained that the propositions incriminated and condemned by the pope as
heretical were not in Jansenius' work entitled \emph{Augustinus}, and he refused
the pope the right to erect into an established fact a point expressly
verifiable by everyone. On an earlier occasion he had already, in his
\emph{Memoir on the Vacuum}, repelled the interference of religious authority in
the problems of the experimental sciences. And in his \emph{18th Provincial
Letter} he sets out that if Galileo's opinion on the movement of the earth is an
error, it is not a decree of the pope that will demonstrate it, and that in the
case where it is possible to prove that it is the earth that turns, it carries
along in its movement those who deny it as well as those who admit it. Pascal
was not a Copernican, but he touches here on a very delicate matter which had
caused Giordano Bruno to be condemned to the fire and had made Galileo perjure
himself. One and the same breath of rebellion against religious authority makes
the secret link between Pascal the physicist and the Pascal of the
\emph{Pensées}. He goes so far as to say that ``the pope hates and fears
scholars who are not bound to him by vow'' [as the religious and the priests
are] (fr. 873, Br.). The remark was aimed at the pope's condemnation of the
Augustinian doctrine of efficacious grace and of the absolute dependence of the
human will in relation to the will of God, after Jansenius, taking it up again,
had accentuated the opposition to scholastic and Jesuit theology. The pope
himself would therefore not be orthodox, and his authority is no more infallible
in matters of faith than when facts are in question. For the Pope has to attend
to so many affairs that he has not the time to occupy himself with the things of
faith
% --- p. 239 ---
(fr. 882). But if Rome condemns the truth, one must cry out all the louder; one
must appeal to heaven. ``It is better to obey God than men'' (fr. 920). Pascal's
Jansenist friends thought it well, subsequently, to suppress these last lines in
the first edition of the \emph{Pensées}. It was a great grief to him to feel
that on this point they did not set truth above everything. Death alone had
prevented him from publishing his last thoughts himself. On his deathbed he
received the sacraments from the hand of his parish priest; his adversaries
interpreted this as a retractation, but one must see in it only that he did not
think he had broken with the true Church.

Kierkegaard grew up in the bosom of the Protestant Church, in veneration of its
leaders, contemporary and former. It was during the years 1847--1850 that he
perceived more and more the disagreement between primitive Christianity and
that of modern times, --- and two of his most penetrating writings, \emph{The
Sickness unto Death} (1849) and \emph{Practice in Christianity} (1850),
announced the great settling of accounts he was going to demand of the Church.
The governing idea of these two works is that men have forgotten, above all
among Protestants, the primitive idea, the prototype for the reconciler. The
dogma of expiation has been made into an acquired datum, and men have settled
without rancour into this imperfect world. Kierkegaard, turning to the
representatives of the Church, exhorts them to recognize what a distance
separates present-day religion from the New Testament, and into what apostasy it
has fallen since the heroic dawn of Christianity. The Church leaders thus
challenged kept silence and even advanced on occasion the pretension of being in
continuity with ``the age of the apostles''; whereupon Kierkegaard launched a
blazing protest, above all in his warlike pamphlet \emph{The Moment} (1855). But
he fell in the thick of the fight and died. He expired without the sacrament,
not wishing to take it from the hand of a pastor. He goes further here in his
revolt against the Church than Pascal did, but he did not reject the sacrament
as such, and on this point he too does not depart from the traditions of the old
Church.

\emph{The Sickness unto Death} is the despair that necessarily arises from the
opposition between the spiritual life imposed on man by Christianity
% ^ The original prints "l'hommé" for "l'homme"; the acute is unmistakable at
%   4x and is a printer's error, kept in the transcription.
and the life assigned him by his natural feelings and appetites. In cases where
a man has not already become conscious of the dilemma that hems him in, he will
still be able to enjoy life as a child would; but from time to time anxiety will
announce itself, putting an end to
% --- p. 240 ---
his security. Then, by weakness, by withdrawal into himself, or by defiance, he
may attempt to lessen his fear and to remain deaf to the Christian demand; but
man will not thus avoid the invisible presence of despair, and will not overcome
it until it has first broken out and invaded everything. In practice Christendom
has grown accustomed --- through unawareness, through the romantic drift of
men's minds, or through simple hypocrisy --- to reducing the great conflict
between the angel and the beast by concessions instead of having done with it by
victories.

In the book he gave under the title \emph{Practice in Christianity},
Kierkegaard lays it down as a principle that, so long as the Saviour has not
returned in his glory, it is for the Christian to take the humiliated Jesus as
his model. We must abstract from what dogma transmits to us about the divine
nature of Jesus and try to make ourselves contemporaries of Jesus of Nazareth,
of the man who said he was the Son of God and who was in consequence insulted
and executed. Instead of that, men have --- like schoolboys --- obtained the
solution of the problem by cheating, sparing themselves the trouble of doing the
calculation for themselves. The parallel imposes itself with the state of soul
of Gethsemane prescribed by Pascal.

More and more, Kierkegaard comes to judge that men have made themselves guilty
of a great betrayal with regard to Christianity: they appropriate the gentleness
and the grace while forgetting what was the necessary condition of them. He
holds that this subreption is above all manifest in Protestantism. Catholicism
could never lower itself to such a platitude. And yet the practice of compromise
goes back to Catholicism. The spirit of Christianity is falsified the day when
following Christ to the letter, while withdrawing from the life of the world,
ceased to be the ordinary rule. The differentiation of Christians into two
distinct categories is a nonsense. But the régime of abuses thus inaugurated was
taken up again by Luther, throwing off the cowl and founding a new Church in
which Christianity and the world would merge in a comfortable arrangement.
``Luther, Luther, Luther, the responsibility that falls on you is great!'' And
Luther had passed through the school of despair. But the subsequent generations
lay claim to the brilliant deeds of their ancestors without showing the least
disposition to accomplish any themselves. Protestantism is a plebeian doctrine
which tends to level the world of the spirit downwards. And in large part it is
family life that is guilty. By
% --- p. 241 ---
means of infant baptism, all the inhabitants of so-called Christian countries
were enrolled in Christendom. ``Men dare,'' cries Kierkegaard (\emph{The
Moment}, VII), ``men dare do this in the face of God, under cover of Christian
baptism! Baptism, the sacred act by which the Saviour of the world was
consecrated to the work of his life, and, after him, the disciples --- men long
since of an age to know, and who, dead to the life of this world, promised to
live henceforth as sacrificial victims in this world of lying and of malice!''

Kierkegaard's polemic had for its culmination the enunciation of a single
proposition (as against Luther's ninety-five): ``The Christianity of the New
Testament does not exist.'' And, completing the thesis, Kierkegaard warned
against ``those who wear long robes'' and against participation in the official
cult, where God is mocked!

The opposition of Pascal and Kierkegaard is a fact of the highest interest in
the history of religion. On the side of the sociologists it has been maintained
that religion is essentially a social phenomenon. And in fact the history of
religions studies by preference religious phenomena belonging to traditional
cult. Under these conditions, how can one attribute an importance of the same
order to figures who were brought, by the effort of their life and their
thought, to break with the religious society to which they belonged, without
afterwards entering any other? If at least they allowed themselves to be
attached each to his sect: Pascal to Jansenism, Kierkegaard, possibly, to
Pietism (as has been done in an excellent French work of sociology). The
attachment would fail, first of all, for Pascal, who precisely, in the supreme
and decisive phase of his life, had the pain of separating from his allies of
Port-Royal; and as for Kierkegaard, he abounds in ironic sallies and in sarcasms
directed at the Pietists and appears in no way inclined to make common cause
with them. It would be more plausible to say that their rupture with the
tradition reigning in the Church of their time had no other aim than to raise up
an older one, in their eyes the only legitimate one. But their appeal, their
return towards what they believed to be the original state of things, made them
solitaries of their epoch, and it was in the name of individual conscience that
they protested. To what social authority could they have
% --- p. 242 ---
appealed? Pascal did, it is true, see in a miracle the confirmation of his
doctrine; only, the miracle in question had received no ecclesiastical
ratification. And Kierkegaard, for his part, expressly affirmed that the type of
religion which alone had value for him no longer existed. In their respective
attitudes towards the sacrament of the last hour a link subsists, too weak
however to bring them back into contact with the Church. --- For the rest, it is
an untenable theory to occupy oneself to such a degree with tradition and with
cult that one comes to neglect the earlier vicissitudes of the souls amenable to
these social forces. One of the dominant tasks of the study of religions must be
to trace step by step the stages of the great religious souls. And this all the
more in that the evolutionary line of these superior individualities traverses
the religion in question to its extreme confines, precisely by their very effort
to integrate its spirit and its truth into their existence. It is just such an
effort that Pascal and Kierkegaard realized in their religious development.

\section*{IV. --- \textsc{Courses to be taken.}}
% ^ The original head is "Partis a prendre.", with the preposition set "A"
%   unaccented, as this journal's capitals always are (confirmed at 400 dpi
%   against "TEMPÉRAMENT"/"CARACTÈRE" on p. 225, which do carry accents).

Conjectures are vain as to the direction the activity of these two great minds
might have taken had their lives been prolonged. But is it without interest to
consider the ways that might open to any man who had followed the two
protagonists, with a conscious adherence reinforced by his own experience, up to
the exact point of their last thought?

Three possibilities present themselves.

One might take up again the principles of primitive Christianity and strive to
make of them the foundations of a life led in intense expectation of God's
approaching intervention, abstaining entirely from collaborating in civilization
and in the current works of human interest. There exist perhaps --- who knows?
--- men engaged on this road, which Saint Francis of Assisi followed in his own
manner.

Another course to be taken would be to begin by making the concession Kierkegaard
demanded, concerning the gap between primitive Christianity and that of our own
day. One would then have to take a further step in this direction and recognize
that at bottom modern Christianity is another religion than primitive
Christianity. One would thus avoid the contradiction that consists in professing
an ideal which in practice one denies. On this conception one would indeed see
in religion an essential factor of the spiritual life, but one would admit
% --- p. 243 ---
alongside it other ideals differing from it in their origin and in their
orientation, which nevertheless deserve to direct in part the conduct and the
conception of life. It will perhaps be difficult to coordinate the different
guiding stars in one and the same system; but at least the sincerity so dear to
our two thinkers will remain intact.

For those to whom these two roads appear impracticable, a third remains open:
they will begin again the effort of the Greeks to found the directives of a moral
life upon human nature and human experience. They will have to seek out the
ideal and the end respectively involved in the principles thus established. In
fact it is in this direction that for two or three centuries the attempts to
realize a morality on a philosophical basis have been moving. I once expressed
this idea (\emph{La Morale}, X, 4) in the following form: ``We can take up again
the Platonic notion of \emph{justice}, as personal harmony, to express the
intimate connection, the higher unity of self-affirmation and of devotion. The
ancient notion of harmony must be enlarged and deepened by means of the moral
experiences that are condensed in the appraisals of character made by Stoicism,
Christianity, the Renaissance and modern Humanity. But this change in the content
of the notion is quite possible without any need to burst its frame. Our morality
is Greek morality, which later moral evolution, with its multiple oscillations,
may well have modified and rectified, but which can never be supplanted so long
as there subsists a human morality truly worthy of the name.'' Continuous work on
these foundations cannot but develop belief in its own future, the conviction
that it has its place in universal existence and even represents an element
indispensable to its full development. For the more living expression of such a
conviction, one must have recourse to the symbolizations of poetic language. And
so poetry, rich and profound, is among the values whose permanence and
flowering in the future are to be counted on.

Over against the two other conceptions of life that we have indicated, the
position to be taken seems to us to be this: instead of renouncing all natural
development of human life and all autonomous culture, family life and all social,
artistic or scientific activity, and instead of attempting a superficial
combination of
% --- p. 244 ---
oriental dogma and occidental culture, one will on the contrary take up
resolutely the cause of human nature and occidental culture and work
courageously for their advancement. Problems to be solved will not be lacking,
nor disappointments; failures will occur. Disputes may arise over the roads to
be taken and over the establishment of guiding ideas, but the general
orientation will always let itself be recognized. Other conceptions of life will
be drawn upon; they will furnish examples, useful nourishment. The lyrical
impulse the believer puts into his expectation of the kingdom of God, always
near, will have a more or less vigorous counterpart in every active aspiration
towards a lofty end. Every work of great scope involves a certain tension. If a
comparison had to be established between the influence Christianity has had on
human life in Europe and that exercised by Buddha in Asia, I should say, to take
up again a formula from my \emph{Philosophie de la religion}, that Buddha
softened Asia and that Christ taught Europe the infinite exaltation of effort.
But above all, the spirit of charity preached by Christianity and placed by it
(with a logic sometimes failing) at the head of the other virtues will have its
appointed place in a morality instituted on purely human foundations. Aristotle
maintained already that justice takes on, in its highest form, the character of
love; the converse is true as well: love is truly a virtue only if, uniting
itself with wisdom, it becomes justice. Generally speaking, the purely human
conception stands in the same relations to Christianity as to the other spiritual
movements: it assimilates what it judges compatible with its continuous
aspiration. As much as the civilizations, the great universalist religions draw
their origins from human nature and from the conditions of man's existence; the
fruits they bear, the sum of nobility, of grandeur, of beauty that we owe them
belongs therefore, in the last instance, to humanity entire and not to a single
sect, however vast its extension may be.

Whatever the road on which one finds oneself engaged, one will draw profit from
the teaching of the two great figures that have just held our attention. Thanks
to them, the centre of gravity of the spiritual life has passed from the outer
world of authorities and dogmas into the inner field of personality, of its
experiences and its needs. Christianity, Pascal concludes, is the true religion
% --- p. 245 ---
because it alone answers to the deepest aspiration of human nature, because it
alone reconciles the greatest contraries of life and gives an account of them.
And for Kierkegaard the idea of God was, in the storms of life, an anchor of
salvation whose necessity does not impose itself on one who keeps always in calm
water. The essential thing, for both, is that each should pass, in the course of
life, through the profitable and penetrating cycle of experience. But first and
above all, one must avoid deceiving oneself and deceiving others by professing
beliefs that are not founded on personal experiences and aspirations. ``In
religion,'' says Pascal (fr. 590 Br.), ``one must be sincere; true pagans, true
Jews, true Christians.'' Kierkegaard, for his part, said during his last campaign
against the established Church: ``I do not dare to claim for myself the name of
Christian, but I range myself under the banner of sincerity; it is for that that
I expose myself.''
% ^ The original closes this second quotation with the guillemet BEFORE the
%   full stop, where the Pascal quotation above closes with the stop inside.
%   The journal's own inconsistency; the transcription leaves it standing.

On one point, however, the men engaged on the third road part company sharply
with Pascal and Kierkegaard. The two thinkers realize that life, and
particularly religious life, needs thoughts to maintain and direct its
orientation. But they make a sharp distinction between vital thoughts and
scientific thoughts or, to speak Kierkegaard's language, between ``essential''
and ``non-essential'' knowledge. Under non-essential knowledge they rank science,
all the sciences. And yet there are scientific thoughts that may become of
capital importance for personal beliefs. Impossible to mark the limits here once
and for all. Let us only point to the influence, muffled and little remarked by
the general public, that is exercised on religious ideas by the admirable
advance the sciences have made --- natural sciences as well as moral sciences
--- during the last centuries. And let us note in this connection how ill founded
is the opinion that would have personality and science stand opposed to one
another as the interests they cover would be opposed. Science has issued from the
work of individuals and is the daughter of personal aspiration; the joy of
knowing counts among the noblest feelings of man. The very frames, conscious or
unconscious, of intellectual life are furnished, in the last analysis, by
personality, since indeed they depend on human nature and the human mind. On the
other hand, personality develops through scientific work; it
% --- p. 246 ---
gains thereby a lucid understanding of life and its conditions, and learns
thereby to impose on itself tasks that far surpass its interested egoisms. And
it is with every other culture as with science. Unceasingly, man is now superior,
now inferior, to this or that given culture. He dominates it because no
discipline justifies its value except by contributing to the development of the
personal life; and conversely culture dominates man by imposing tasks that
imperiously subordinate individual interests to the lasting gain, the lasting
progress, of the human collectivity.

\bigskip
% ^ The original sets a blank line here --- more than a paragraph break, less
%   than a section break --- before the closing paragraph, as on p. 226.

Two great figures, two solitaries without continuators or disciples. But the
analogy of their governing ideas and of their respective attitude in the face of
the realities of their time contain a lesson for every epoch.
% ^ The original prints "contiennent", plural, against the singular subject
%   "l'analogie". Høffding's own French; the plural is carried over here.
From their high solitude they discerned great and decisive problems in the field
of the history of religions and of religious psychology, and they left their
lives in their energy to pose these problems in an unforgettable way. The more
unforgettable in that they poured into them not only passion but such a wealth
of ideas and so dense a poetry that these rank their works among the most
precious of human treasures.

\bigskip

\hfill \textsc{Harald Höffding.}

\end{document}
